\documentclass{article}

% Packages
\usepackage[final]{neurips_2024}
\usepackage[utf8]{inputenc}
\usepackage[T1]{fontenc}
\usepackage{hyperref}
\usepackage{url}
\usepackage{booktabs}
\usepackage{amsfonts}
\usepackage{amsmath}
\usepackage{amssymb}
\usepackage{microtype}
\usepackage{graphicx}
\usepackage{multirow}
\usepackage{xcolor}
\usepackage{algorithm}
\usepackage{algorithmic}
\usepackage{footnote}

\title{Multi-Constraint Quantum Document Selection for\\
Retrieval-Augmented Generation:\\
An Empirical Study of Classical, Annealing,\\
and Gate-Based Quantum Solvers}

\author{
  Joe Reinert \\
  Independent Researcher \\
  Cincinnati, OH \\
  \texttt{jreinert11@outlook.com} \\
}

\begin{document}

\maketitle

% ─────────────────────────────────────────────
\begin{abstract}
Retrieval-Augmented Generation (RAG) systems typically select context
documents using greedy nearest-neighbor search, which optimizes each
document independently and ignores combinatorial interactions between
retrieved items. We present \textbf{TaxQRAG}, the first open-source
quantum-hybrid RAG system evaluated on real quantum hardware for
multi-strategy document selection under explicit pairwise interaction
constraints. Unlike prior work, our QUBO (Quadratic Unconstrained
Binary Optimization) formulation incorporates three novel quadratic
terms: domain-specific pairwise \emph{synergy rewards}, \emph{conflict
penalties}, and \emph{audit risk compounding} — enabling selection of
document subsets that are jointly optimal rather than individually
relevant. Applied to a tax strategy corpus (87 documents; IRS
publications, Internal Revenue Code sections, and structured strategy
documents), we compare four solvers: classical greedy, simulated
annealing (SA), QAOA on a local simulator (Aer), and QAOA on IBM
Quantum's Heron R2 processor (\texttt{ibm\_fez}, 156 qubits). SA
achieves a 173\% improvement in combination quality score over classical
greedy with zero conflict selections. QAOA without parameter
optimization produces a 66.7\% conflict rate, which a $\lambda$
ablation study confirms is caused by unoptimized variational parameters
rather than QUBO encoding deficiency — a negative result that
motivates VQE-optimized QAOA as future work. We further conduct the
first QAOA parameter sensitivity analysis for RAG document selection,
finding that combination quality coefficient of variation reaches 88.2\%
across random seeds, and that Ising energy and downstream combination
score diverge significantly at shallow circuit depth. Code, benchmark
data, and IBM Quantum job identifiers are released at
\url{https://github.com/joerein-iu/taxqrag}.
\end{abstract}

% ─────────────────────────────────────────────
\section{Introduction}

Retrieval-Augmented Generation \cite{lewis2020rag} has emerged as the
dominant paradigm for grounding large language model (LLM) responses in
external knowledge. The standard retrieval pipeline embeds a query,
retrieves the $K$ nearest documents by cosine similarity, and passes
them to the LLM as context. This approach has a fundamental limitation:
it treats document selection as $K$ independent scalar comparisons
rather than a combinatorial optimization over document subsets.

In domains where documents interact — where strategy A amplifies
strategy B, or strategy C legally conflicts with strategy D — greedy
nearest-neighbor selection is provably suboptimal. The combinatorial
problem of selecting the subset of $K$ documents that maximizes joint
value subject to multiple constraints is NP-hard classically, making it
a natural candidate for quantum optimization.

Prior work on quantum-enhanced RAG \cite{qrag2025} has demonstrated
Grover's algorithm for query acceleration and QAOA for basic relevance
ranking, but has been limited to: (1) single-objective QUBO
formulations, (2) simulated quantum hardware, and (3) no open-source
implementations for reproducibility.

\textbf{This paper makes four contributions:}
\begin{enumerate}
    \item A novel multi-constraint QUBO formulation incorporating
    synergy rewards, conflict penalties, and audit risk compounding as
    quadratic interaction terms (Section~\ref{sec:method}).
    \item The first open-source implementation evaluated on real IBM
    Quantum hardware for RAG document selection (Section~\ref{sec:experiments}).
    \item A rigorous three-way empirical comparison of classical greedy,
    simulated annealing, and gate-based QAOA solvers with honest
    reporting of failure modes (Section~\ref{sec:results}).
    \item A QAOA parameter sensitivity analysis and $\lambda$ ablation
    study establishing that QAOA's conflict rate stems from
    unoptimized variational parameters, not QUBO formulation
    (Section~\ref{sec:analysis}).
\end{enumerate}

% ─────────────────────────────────────────────
\section{Background}
\label{sec:background}

\subsection{Retrieval-Augmented Generation}

RAG systems \cite{lewis2020rag} retrieve relevant documents from an
external corpus and condition LLM generation on the retrieved context.
Standard retrieval uses dense vector search (e.g., HNSW \cite{hnsw})
to find nearest neighbors in embedding space. Reranking approaches such
as Maximal Marginal Relevance (MMR) \cite{carbonell1998mmr} partially
address redundancy but optimize greedily and cannot encode arbitrary
pairwise constraints.

\subsection{QUBO and Quantum Annealing}

A Quadratic Unconstrained Binary Optimization (QUBO) problem is defined
as:
\begin{equation}
\min_{x \in \{0,1\}^n} x^T Q x
\end{equation}
where $Q \in \mathbb{R}^{n \times n}$ encodes both linear (diagonal)
and quadratic (off-diagonal) objective terms. QUBO is NP-hard in
general but is the native problem format for quantum annealers
(D-Wave) and can be transformed to the Ising Hamiltonian for gate-based
QAOA via the substitution $x_i = (1 - z_i)/2$.

\subsection{QAOA}

The Quantum Approximate Optimization Algorithm \cite{farhi2014qaoa}
applies alternating problem and mixer layers to a quantum circuit,
parameterized by angles $\boldsymbol{\gamma}, \boldsymbol{\beta}$:
\begin{equation}
|\psi(\boldsymbol{\gamma}, \boldsymbol{\beta})\rangle =
\prod_{l=1}^{p} e^{-i\beta_l H_M} e^{-i\gamma_l H_C} |+\rangle^n
\end{equation}
where $H_C$ is the cost Hamiltonian encoding the QUBO objective and
$H_M$ is the transverse-field mixer. Optimal parameters are found via
classical optimization (VQE); without this step, random initialization
samples an arbitrary point in the variational landscape.

% ─────────────────────────────────────────────
\section{Method}
\label{sec:method}

\subsection{Problem Formulation}

Given a query $q$ and corpus $\mathcal{D}$, we seek to select a subset
$S \subseteq \mathcal{D}$, $|S| = K$ that maximizes joint combination
quality subject to a token budget $B$. Let $x_i \in \{0,1\}$ indicate
selection of document $i$. Our QUBO objective is:

\begin{equation}
\min_{x} \; \underbrace{-\sum_i r_i x_i}_{\text{relevance}} \;
+ \underbrace{\sum_{i < j} C_{ij} x_i x_j}_{\text{conflict}} \;
- \underbrace{\sum_{i < j} S_{ij} x_i x_j}_{\text{synergy}} \;
+ \underbrace{\lambda_{\text{tok}} \left(\sum_i t_i x_i - B\right)^2}_{\text{token budget}} \;
+ \underbrace{\lambda_{\text{cnt}} \left(\sum_i x_i - K\right)^2}_{\text{cardinality}}
\label{eq:qubo}
\end{equation}

where $r_i = \alpha \cdot \text{rel}_i + \beta \cdot \text{cred}_i$ is
a weighted linear relevance term combining cosine similarity
$\text{rel}_i$ and source credibility $\text{cred}_i \in [0, 1]$
(IRS publications: 1.0; IRC sections: 1.0; structured strategies: 0.95).

\subsection{Interaction Matrices — The Novel Contribution}

Prior QUBO RAG work \cite{qrag2025} uses semantic cosine similarity for
quadratic terms, penalizing redundant document pairs. We replace this
with \emph{domain-specific interaction matrices} derived from tax law:

\textbf{Synergy matrix} $S_{ij} \in [0, 1]$: reward for selecting
complementary strategy pairs. Example: S-Corporation salary
optimization and Solo 401(k) contributions have synergy $S_{ij} = 0.95$
because the W-2 salary determines the maximum employer match —
selecting both is more valuable than selecting either alone.

\textbf{Conflict matrix} $C_{ij} \in [0, 1]$: penalty for selecting
legally conflicting strategies. Example: the Augusta Rule (IRC §280A(g))
and home office deduction conflict at $C_{ij} = 0.95$ because the same
physical space cannot simultaneously qualify for both treatments under
IRS rules.

\textbf{Audit risk compounding} $A_{ij} \in [1, 2]$: multiplicative
audit risk escalation for certain strategy combinations. Example:
Augusta Rule combined with S-Corporation salary optimization compounds
audit risk at $A_{ij} = 1.4\times$, encoded as an additional quadratic
penalty term $\lambda_{\text{aud}}(A_{ij} - 1)$.

Table~\ref{tab:interactions} shows a subset of the 20 encoded strategy
pairs. The full interaction matrix is released with the codebase.

\begin{table}[h]
\centering
\caption{Selected pairwise interaction scores. Synergy rewards joint
selection; conflict penalizes it. Audit multiplier compounds risk.}
\label{tab:interactions}
\small
\begin{tabular}{llccc}
\toprule
Strategy A & Strategy B & Synergy & Conflict & Audit \\
\midrule
S-Corp Salary & Solo 401(k) & 0.95 & 0.00 & 1.00× \\
S-Corp Salary & QBI Deduction & 0.85 & 0.00 & 1.00× \\
S-Corp Salary & Augusta Rule & 0.80 & 0.00 & 1.40× \\
TTS Election & §1256 Contracts & 0.90 & 0.00 & 1.00× \\
TTS Election & Solo 401(k) & 0.75 & 0.00 & 1.00× \\
Augusta Rule & Home Office & 0.00 & 0.95 & 2.00× \\
Solo 401(k) & SEP-IRA & 0.00 & 0.80 & 1.00× \\
§1256 Contracts & S-Corp Trading & 0.00 & 0.85 & 1.00× \\
\bottomrule
\end{tabular}
\end{table}

\subsection{Corpus}

We construct a 87-document tax strategy corpus from public domain
sources: 10 hand-crafted strategy documents encoding structured
knowledge about each strategy (qualifications, IRC authority, synergy
metadata), 24 IRC section documents fetched from Cornell LII, and 53
chunked IRS publication segments from Publications 535, 550, 560, 590a,
969, and 946. All sources are U.S. government works in the public
domain. Documents are embedded using \texttt{all-MiniLM-L6-v2}
\cite{reimers2019sbert} and stored in ChromaDB for dense retrieval.

\subsection{Pipeline}

At query time: (1) embed query and retrieve $K_c = 30$ candidates via
cosine similarity; (2) formulate QUBO over candidates using
Equation~\ref{eq:qubo}; (3) solve via one of three backends; (4) decode
solution to selected document subset; (5) pass subset to Claude
(claude-sonnet-4-20250514) for generation.

% ─────────────────────────────────────────────
\section{Experiments}
\label{sec:experiments}

\subsection{Solvers}

We evaluate three solvers against the same QUBO objective:

\textbf{Classical Greedy:} Sort candidates by relevance score; select
greedily subject to token budget and hard conflict avoidance (skip
documents with conflict score $> 0.8$ against any already-selected
document). This represents the industry baseline.

\textbf{Simulated Annealing (SA):} \texttt{dimod.SimulatedAnnealingSampler}
with \texttt{num\_reads=50}, \texttt{num\_sweeps=500}. Runs locally with
no cloud dependency.

\textbf{QAOA (Aer Simulator):} \texttt{QAOAAnsatz} at $p=2$ (two
alternating layers), local Aer simulation, 1024 shots, seed=43,
\texttt{optimization\_level=1} transpilation.

\textbf{QAOA (IBM Quantum):} Same QAOA circuit submitted to
\texttt{ibm\_fez} (IBM Heron R2, 156 qubits) via IBM Quantum Open Plan.
20 qubits used (top-20 candidates by relevance). Single query result;
job ID \texttt{d7i5h23jne2c7394sh60}.

\subsection{Benchmark}

12 tax strategy queries spanning self-employed trader scenarios,
small business owner scenarios, and explicit conflict-testing scenarios
(e.g., ``Can I use the Augusta Rule and home office deduction at the
same time?''). Classical and SA results are averaged over 2 runs per
query (24 total evaluations each). QAOA Aer results use a fixed seed
for reproducibility. IBM QPU result is a single query run due to free
tier constraints (10 minutes per 28-day rolling window).

\subsection{Metrics}

\textbf{Combination Quality Score:} $\sum_{i<j, i,j \in S} S_{ij} -
C_{ij} \cdot 2 - (A_{ij} - 1)$ for selected set $S$. Captures the
net value of the selected combination accounting for synergies,
conflicts, and audit risk.

\textbf{Total Synergy Captured:} $\sum_{i<j, i,j \in S} S_{ij}$.
Raw synergy without conflict penalty.

\textbf{Conflict Rate:} Fraction of queries where selected set contains
at least one conflicting strategy pair ($C_{ij} > 0.7$).

\textbf{Average Relevance:} Mean cosine similarity of selected
documents to query embedding.

% ─────────────────────────────────────────────
\section{Results}
\label{sec:results}

\subsection{Main Benchmark}

Table~\ref{tab:main} reports aggregate results across all solvers.

\begin{table}[h]
\centering
\caption{Main benchmark results. QAOA Aer uses seed=43, $p=2$,
random-init parameters. IBM QPU: single-query result on \texttt{ibm\_fez}
(Heron R2). $^\dagger$High variance; see Section~\ref{sec:analysis}.
$^\ddagger$Single query, $n=1$.}
\label{tab:main}
\begin{tabular}{lcccc}
\toprule
Metric & Classical & Sim. Anneal & QAOA (Aer)$^\dagger$ & IBM QPU$^\ddagger$ \\
\midrule
Combination Score & 0.675 & 1.846 & 2.779 & 5.100 \\
Total Synergy & 0.700 & 1.846 & 4.371 & 5.500 \\
Synergies Found & 0.75 & 1.875 & 4.583 & 5.0 \\
Conflict Rate & 0.0\% & 0.0\% & 66.7\% & 0.0\% \\
Avg. Relevance & 0.447 & 0.377 & 0.362 & — \\
Rerank Time & $<$1ms & 7.4s & 605ms & 21s \\
\midrule
$\Delta$ vs. Classical & — & +173.5\% & +311.7\%$^\dagger$ & +655\%$^\ddagger$ \\
\bottomrule
\end{tabular}
\end{table}

Simulated annealing achieves the best \emph{practical} result: 173.5\%
improvement in combination quality over classical greedy with zero
conflict selections. QAOA in both configurations captures more raw
synergy than SA but at the cost of a 66.7\% conflict rate (Aer) and
high variance (IBM QPU single-query). Classical greedy is the only
solver with zero conflicts by construction due to its explicit conflict
avoidance heuristic.

The IBM QPU single-query result (combination score 5.100, 5 synergy
pairs) is the highest observed value across all experiments, but
insufficient for statistical claims. We report it for reproducibility
and as evidence of the hardware's ability to find high-quality
solutions. The selected strategies for this run were: \textsc{Section
1256}, \textsc{QBI Deduction}, \textsc{S-Corp Salary}, \textsc{Solo
401(k)}, \textsc{HSA}, \textsc{Augusta Rule}, \textsc{Roth Conversion},
plus supporting IRS publication chunks.

\subsection{Conflict Rate Decomposition}

The 66.7\% conflict rate of QAOA Aer is not distributed uniformly.
Queries explicitly testing conflict scenarios (Q5: Augusta Rule +
home office; Q7: same) both produce combination scores of $-2.9$,
the worst observed values. This indicates that QAOA without parameter
optimization cannot enforce the conflict penalty even when the
conflicting pair is the dominant feature of the query. This observation
motivates our ablation study in Section~\ref{sec:analysis}.

% ─────────────────────────────────────────────
\section{Analysis}
\label{sec:analysis}

\subsection{QAOA Parameter Sensitivity}

We evaluate QAOA's sensitivity to random parameter initialization by
running the same query (``S-Corp active futures trading on MNQ
contracts'') with four seeds (42–45) at optimization level 1 on
\texttt{ibm\_fez}. Table~\ref{tab:sensitivity} reports results.

\begin{table}[h]
\centering
\caption{QAOA parameter sensitivity on \texttt{ibm\_fez}.
All runs: 20 qubits, 1024 shots, $p=1$, random-init parameters.
Job IDs are provided for reproducibility.}
\label{tab:sensitivity}
\small
\begin{tabular}{lccccc}
\toprule
Seed & Depth & Energy & Combo Score & Synergy & Job ID \\
\midrule
42 & 1104 & $-44.39$ & 1.200 & 1.50 & \texttt{d7i5kk493s0c738tjto0} \\
43 & 769  & $+24.63$ & \textbf{6.650} & 7.05 & \texttt{d7i5ktk93s0c738tju70} \\
44 & 972  & $-28.66$ & 1.450 & 1.45 & \texttt{d7i5l2a2khts739qg0eg} \\
45 & 1115 & $+21.46$ & 2.250 & 2.25 & \texttt{d7i5lbvb91ec73avfjq0} \\
\midrule
Mean & 990 & $-6.74$ & 2.888 & 3.06 & — \\
Std & 161 & 35.01 & 2.548 & 2.65 & — \\
CV & 16.3\% & n/a & \textbf{88.2\%} & — & — \\
\bottomrule
\end{tabular}
\end{table}

The 88.2\% coefficient of variation on combination score demonstrates
that random parameter initialization produces highly unstable solutions
on real NISQ hardware. The standard deviation of Ising energy (35.01)
spans both positive and negative values across seeds, indicating that
random angles sample widely separated basins of the variational
landscape.

\subsection{Energy-Combination Score Divergence}

A particularly striking finding emerges from Table~\ref{tab:sensitivity}:
\emph{Ising energy and downstream combination quality are not
monotonically related.} Seed 43, which achieves the worst Ising energy
(+24.63, indicating a high-energy state far from the ground state of
the cost Hamiltonian), simultaneously achieves the best combination
score (6.650). Seed 42, with the best energy at optimization level 3
($-73.61$), achieves the worst combination score (0.700).

This divergence arises because the QAOA objective — minimizing the
Ising Hamiltonian — is only a proxy for the downstream combination
quality metric. At shallow depth ($p=1,2$) without parameter
optimization, the most-frequently sampled bitstring reflects the
structure of the variational ansatz rather than the true ground state.
The ground state of the Ising Hamiltonian would correspond to the
globally optimal document selection, but QAOA without VQE does not
reliably reach it. This misalignment motivates either deeper circuits
($p \geq 3$) or VQE parameter optimization.

\subsection{Transpiler Optimization}

With QAOA parameters held fixed at seed=42, we compare
\texttt{optimization\_level=1} versus \texttt{optimization\_level=3}
transpilation. Level 3 achieves a 39.3\% reduction in circuit depth
(1104 $\to$ 670 gates) and a lower Ising energy ($-44.39 \to -73.61$).
However, the combination score decreases (1.200 $\to$ 0.700),
consistent with the energy-score divergence finding. We conclude that
transpiler optimization improves the fidelity of circuit execution but
does not resolve the parameter optimization problem. Level 3 is
recommended for final paper runs.

\subsection{$\lambda_{\text{cnt}}$ Ablation — Negative Result}

We test whether increasing the cardinality penalty weight from
$\lambda_{\text{cnt}} = 1.0$ to $5.0$ reduces QAOA's 66.7\% conflict
rate. Table~\ref{tab:ablation} reports results on four
conflict-producing queries.

\begin{table}[h]
\centering
\caption{$\lambda_{\text{cnt}}$ ablation on conflict queries.
QAOA Aer, seed=43, $p=2$. Three of four queries produce
\emph{identical bitstrings} despite $4\times$ energy scaling.}
\label{tab:ablation}
\small
\begin{tabular}{lcccc}
\toprule
Metric & $\lambda=1.0$ & $\lambda=5.0$ & $\Delta$ \\
\midrule
Docs selected (mean) & 20.0 & 15.0 & $-5.0$ \\
Combination score & $-1.075$ & $-0.350$ & $+0.725$ \\
Total conflict & 0.712 & 0.475 & $-0.237$ \\
Conflict rate & 100\% & 75\% & $-25$pp \\
Identical bitstrings & — & 3/4 queries & — \\
\bottomrule
\end{tabular}
\end{table}

The key finding: three of four queries produce \emph{bit-identical
output} at $\lambda_{\text{cnt}} = 5.0$ despite the QUBO energy of
those bitstrings increasing by $4\times$ the per-variable penalty delta.
Query 7 is the sole exception, but resolves to the all-zeros bitstring
(zero documents selected) — a degenerate solution that is
trivially conflict-free but useless for RAG.

This result is diagnostic: \emph{QAOA without parameter optimization is
not responsive to the QUBO objective landscape.} Increasing the penalty
raises the energy of bad solutions without changing which solutions are
sampled, because the random variational parameters are not driven toward
the energy minimum. The 66.7\% conflict rate is therefore not a QUBO
encoding deficiency but a fundamental limitation of un-optimized QAOA
on NISQ hardware. The correct fix is VQE parameter optimization, not
QUBO weight tuning.

% ─────────────────────────────────────────────
\section{Discussion and Future Work}
\label{sec:discussion}

\subsection{Practical Recommendation}

For practitioners seeking to deploy quantum-hybrid RAG today, simulated
annealing is the clear recommendation: it achieves 173\% better
combination quality than classical greedy, enforces all constraints
reliably, runs locally without cloud dependency, and completes in under
10 seconds per query. QAOA on current NISQ hardware requires VQE
parameter optimization before it can be used reliably in production.

\subsection{Limitations}

\textbf{QAOA parameter optimization.} All QAOA results in this paper
use random parameter initialization. VQE-optimized parameters would
likely improve solution quality and reduce conflict rate substantially.
We leave this as the primary future work item.

\textbf{IBM QPU statistical power.} The IBM QPU result is a single
query run ($n=1$) due to free tier constraints. The impressive
combination score (5.100) should be interpreted as a proof of concept,
not a statistically robust finding. Multi-query evaluation on real
hardware requires paid access (IBM Quantum Flex or Premium Plan).

\textbf{Corpus size.} The 87-document corpus is small by RAG standards.
Larger corpora (thousands of documents) would require a two-stage
approach: dense retrieval to a larger candidate pool, then QUBO over
the top-$N$ candidates.

\textbf{Interaction matrix coverage.} The current matrix encodes 20
strategy pairs. A comprehensive tax strategy interaction matrix would
require collaboration with licensed tax professionals and coverage of
hundreds of strategy combinations.

\subsection{Future Work}

\textbf{VQE-optimized QAOA.} Warm-starting QAOA parameters from the SA
solution \cite{warmstart2021} is the most promising near-term
improvement. SA finds a good classical solution; QAOA can then be
initialized near that solution and fine-tuned with a small number of
classical optimization steps.

\textbf{Reps scaling.} Increasing QAOA depth from $p=1$ to $p \geq 3$
should improve approximation ratio. The practical limit is circuit depth
vs.\ decoherence on current hardware.

\textbf{Domain generalization.} The interaction matrix approach
generalizes to any domain with explicit pairwise document constraints:
medical treatment combinations, legal precedent compatibility,
investment portfolio diversification. Each domain requires a
domain-specific interaction matrix but the QUBO formulation is
identical.

\textbf{Fault-tolerant hardware.} As quantum hardware matures toward
fault tolerance, the circuit depth constraints that limit current QAOA
performance will relax. The QUBO formulation presented here is hardware-
agnostic and will benefit directly from improved processors.

% ─────────────────────────────────────────────
\section{Related Work}
\label{sec:related}

\textbf{Quantum-enhanced RAG.} \citet{qrag2025} propose QRAG, using
Grover's algorithm for query acceleration and QAOA for basic relevance
ranking. Their work runs on simulators only, uses a single-objective
QUBO (relevance alone), and is not open-sourced. Our work extends this
with multi-constraint QUBO formulation, real hardware evaluation, and
rigorous failure mode analysis.

\textbf{Classical RAG reranking.} Maximal Marginal Relevance (MMR)
\cite{carbonell1998mmr} penalizes redundancy in reranking but optimizes
greedily and cannot encode arbitrary pairwise constraints.
Cross-encoder reranking \cite{nogueira2019passage} improves relevance
scoring but also selects documents independently. Our QUBO approach is
the first to encode domain-specific synergy and conflict as first-class
optimization objectives.

\textbf{QAOA for combinatorial optimization.} \citet{farhi2014qaoa}
introduced QAOA for MaxCut; subsequent work has applied it to portfolio
optimization \cite{mugel2022portfolio} and scheduling
\cite{venturelli2016scheduling}. Our work is the first application of
QAOA to document selection in RAG systems.

\textbf{Parameter sensitivity in QAOA.} \citet{shaydulin2023qaoa}
analyze QAOA parameter landscapes on NISQ hardware, finding high
variance with random initialization — consistent with our findings.
\citet{warmstart2021} propose warm-starting QAOA from classical
solutions, which we identify as key future work.

% ─────────────────────────────────────────────
\section{Conclusion}
\label{sec:conclusion}

We presented TaxQRAG, a quantum-hybrid RAG system that formulates
document selection as a multi-constraint QUBO incorporating pairwise
synergy rewards, conflict penalties, and audit risk compounding.
Evaluated on a tax strategy corpus using IBM Quantum's Heron R2
processor, we find that simulated annealing achieves 173\% better
combination quality than classical greedy with zero conflicts —
making it the practical recommendation for deployment today. Gate-based
QAOA without parameter optimization produces a 66.7\% conflict rate
that a $\lambda$ ablation study confirms is caused by unoptimized
variational parameters, not QUBO formulation deficiency. This negative
result precisely motivates VQE-optimized QAOA as future work. All code,
benchmark data, and IBM Quantum job identifiers are released for
reproducibility at \url{https://github.com/joerein-iu/taxqrag}.

% ─────────────────────────────────────────────
\section*{Acknowledgments}

We acknowledge the use of IBM Quantum services for this work. The views
expressed are those of the authors, and do not reflect the official
policy or position of IBM or the IBM Quantum team.

IBM Quantum hardware used: \texttt{ibm\_fez} (Heron R2, 156 qubits).
Job identifiers: \texttt{d7i5h23jne2c7394sh60} (main result, opt.
level 1), \texttt{d7i5idjjne2c7394sihg} (opt. level 3),
\texttt{d7i5kk493s0c738tjto0}–\texttt{d7i5lbvb91ec73avfjq0}
(sensitivity analysis, seeds 42–45).

% ─────────────────────────────────────────────
\bibliographystyle{plainnat}
\begin{thebibliography}{99}

\bibitem{lewis2020rag}
Lewis, P., Perez, E., Piktus, A., Petroni, F., Karpukhin, V., Goyal,
N., ... \& Kiela, D. (2020). Retrieval-Augmented Generation for
Knowledge-Intensive NLP Tasks. \textit{Advances in Neural Information
Processing Systems}, 33, 9459--9474.

\bibitem{qrag2025}
Adithi, K., \& Kapilavani, R. K. (2025).
Quantum Synergy in Retrieval-Augmented Generation for Contextual
Enhancement. \textit{Preprint, Research Square}.
DOI: 10.21203/rs.3.rs-6216441/v1.

\bibitem{hnsw}
Malkov, Y. A., \& Yashunin, D. A. (2018). Efficient and Robust
Approximate Nearest Neighbor Search Using Hierarchical Navigable Small
World Graphs. \textit{IEEE Transactions on Pattern Analysis and Machine
Intelligence}, 42(4), 824--836.

\bibitem{carbonell1998mmr}
Carbonell, J., \& Goldstein, J. (1998). The Use of MMR,
Diversity-Based Reranking for Reordering Documents and Producing
Summaries. \textit{Proceedings of SIGIR}, 335--336.

\bibitem{nogueira2019passage}
Nogueira, R., \& Cho, K. (2019). Passage Re-ranking with BERT.
\textit{arXiv preprint arXiv:1901.04085}.

\bibitem{farhi2014qaoa}
Farhi, E., Goldstone, J., \& Gutmann, S. (2014). A Quantum Approximate
Optimization Algorithm. \textit{arXiv preprint arXiv:1411.4028}.

\bibitem{reimers2019sbert}
Reimers, N., \& Gurevych, I. (2019). Sentence-BERT: Sentence Embeddings
using Siamese BERT-Networks. \textit{Proceedings of EMNLP}.

\bibitem{warmstart2021}
Egger, D. J., Marecek, J., \& Woerner, S. (2021). Warm-Starting Quantum
Optimization. \textit{Quantum}, 5, 479.

\bibitem{shaydulin2023qaoa}
Shaydulin, R., Lotshaw, P. C., Larson, J., Ostrowski, J., \& Humble,
T. S. (2023). Parameter Transfer for Quantum Approximate Optimization
of Weighted MaxCut. \textit{ACM Transactions on Quantum Computing}, 4(3).

\bibitem{mugel2022portfolio}
Mugel, S., Kuchkovsky, C., Sanchez, E., Fernandez-Lorenzo, S.,
Luis-Hita, J., Lizaso, E., \& Orus, R. (2022). Dynamic Portfolio
Optimization with Real Datasets Using Quantum Processors and Quantum-
Inspired Tensor Networks. \textit{Physical Review Research}, 4(1).

\bibitem{venturelli2016scheduling}
Venturelli, D., Marchand, D. J., \& Rojo, G. (2016). Quantum Annealing
Implementation of Job-Shop Scheduling. \textit{arXiv preprint
arXiv:1506.08479}.

\end{thebibliography}

\end{document}
